% Propietario reproducido por unidad demostrativa; véase PROPIETARIOS_CONTINUO_UNIVERSO.json.
% Origen: XI ES CUMPLIMIENTO_EDITORIAL_20260928, sections/continuo_cinco_construcciones_tipado.tex:1--327.
% Cuerpos y pruebas conservados; sólo namespace, rutas y remisiones contextuales declaradas.
\subsection{Émergence conjointe des cinq constructions}
\label{hmtfg:base:sec:cinco-construcciones-continuo-ch}

Pour \(0\leq k\leq N\), soit
\[
 \rho_{N,k}:=\rho_k\circ\rho_{k+1}\circ\cdots\circ\rho_{N-1},
 \qquad \rho_{k,k}:=\operatorname{id}_{H_k},
\]
et, pour une histoire survivante \(\zeta\in H_k\), définissons
\[
 \operatorname{Desc}_{k,N}(\zeta)
 :=\{x\in H_N:\rho_{N,k}(x)=\zeta\}.
\]
L'arbre fini \(\mathscr T_{k,N}(\zeta)\) a pour sommets
\[
 \coprod_{j=k}^{N}\{y\in H_j:\rho_{j,k}(y)=\zeta\}
\]
et pour arêtes uniquement les prolongements TPK \(a:y\to z\) avec \(z\in\rho_j^{-1}(y)\). Ainsi, ni ses sommets ni ses prolongements ne proviennent d'un arbre auxiliaire.

Pour \(x\in\operatorname{Desc}_{k,N}(\zeta)\), la branche marquée \(\lambda_{\zeta,x}\) est la suite
\[
 \zeta=x_k\longrightarrow x_{k+1}\longrightarrow\cdots
 \longrightarrow x_N=x
\]
déterminée par les troncatures de \(x\). Le registre de base commun est
\[
\begin{aligned}
 \mathscr R_{k,N}(\zeta;x):=
 \bigl(&\mathscr T_{k,N}(\zeta),\lambda_{\zeta,x};\\
 &\bigl(\varepsilon_a,s_a,o_a,
   (\delta_{\diamond,a},r_{\diamond,a},q_{\diamond,a})
       _{\diamond\in\{\Sigma,\Pi\}},
   \operatorname{Carr}_a,C_a,t_a,\operatorname{Fr}_a,
   \operatorname{Ruta}_a,\operatorname{Mem}_a,
   \operatorname{Surv}_a\bigr)_a;\\
 &\operatorname{src},\operatorname{tgt},\circ,
   (\rho_j)_{k\leq j<N}\bigr),
\end{aligned}
\tag{C1}
\label{hmtfg:base:eq:base-comun-registros}
\]
où
\[
 U_a(m_0)=C_am_0+t_a,
 \qquad
 \delta_{\diamond,a}=r_{\diamond,a}+9q_{\diamond,a}.
\]
Le symbole \(m_0\) de cette formule est une coordonnée de mémoire ; ce n'est pas la profondeur modulaire \(m\) utilisée plus bas. Le véritable espace de base est
\[
 \mathfrak R_{k,N}(\zeta)
 :=\{\mathscr R_{k,N}(\zeta;x):
       x\in\operatorname{Desc}_{k,N}(\zeta)\}.
\]
La branche marquée n'est pas une entrée ajoutée : c'est la suite des arêtes
effectivement produites qui se termine en \(x\). Le reste du registre conserve
simultanément les fibres latérales de l'arbre ; la marque ne remplace donc pas
la ramification par une trajectoire isolée.

\Needspace{15\baselineskip}
Posons \(R_m:=\mathbb Z/3^m\mathbb Z\), et soit
\(V_j(\zeta):=\operatorname{Desc}_{k,j}(\zeta)\). Le transport, les
projecteurs de feuillet et le résidu terminal se restreignent au sous-arbre par
\[
 \mathcal U_j^\zeta:=
 \mathcal U_j\!\upharpoonright_{\ell^2(V_j(\zeta),\mu_j)},
 \qquad P_j^\zeta:=P_j\!\upharpoonright_{V_j(\zeta)},
 \qquad Q_j^\zeta:=I-P_j^\zeta,
\]
\[
 A_j^\zeta:=P_{j+1}^\zeta\mathcal U_j^\zeta P_j^\zeta
             +Q_{j+1}^\zeta\mathcal U_j^\zeta Q_j^\zeta,
 \qquad
 \eta_j^\zeta:=Q_{j+1}^\zeta\mathcal U_j^\zeta P_j^\zeta
             +P_{j+1}^\zeta\mathcal U_j^\zeta Q_j^\zeta.
\]
Pour \(k\leq n\leq N\), soient
\[
 F_{k,k}^\zeta:=I,
 \qquad
 F_{k,n}^\zeta:=A_{n-1}^\zeta\cdots A_k^\zeta,
 \qquad
 T_{k,n}^\zeta:=(F_{k,n}^\zeta)^*F_{k,n}^\zeta.
\]
Les ports propres du même sous-arbre sont
\[
 \mathscr P_j^\Sigma(\zeta):=
 \{(y,\Sigma):y\in V_j(\zeta)\},
 \qquad
 \mathscr P_j^\Pi(\zeta):=
 \{(y,\Pi):y\in V_j(\zeta)\};
\]
nous désignons par \(B_j^\zeta,\kappa_j^\zeta,\delta_j^\zeta\) les restrictions de l'incidence et du complexe rectangulaire à ces ports, et par \(L_j^\zeta\) la précomposition induite par la suppression de préfixes.

Pour définir sans abréviation le lecteur de chemins, soit
\[
 W_\ell^\zeta(y):=
 R_m[\operatorname{Path}_{\mathscr T_{k,\ell}(\zeta)}
          (y,\operatorname{Term}_\ell)]
\]
et posons
\[
 \mathscr D_{q,r;k,\ell}^{(m)}[\mathscr R]
 :=\bigoplus_{y_0\to\cdots\to y_q\ \text{dans}\
                  \mathscr T_{k,\ell}(\zeta)}
 W_\ell^\zeta(y_q)\otimes_{R_m}
 \bigl(\Gamma_r(y_0)\otimes_{\mathbb Z}R_m\bigr),
 \qquad
 \operatorname{Tot}_d\mathscr D
 :=\bigoplus_{q+r=d}\mathscr D_{q,r;k,\ell}^{(m)}[\mathscr R].
\]
Les formules déjà démontrées produisent, sur chaque
\(\mathscr R\in\mathfrak R_{k,N}(\zeta)\), les cinq objets
\[
\begin{aligned}
 \mathcal F_{\rm sp}[\mathscr R]
 &:=(\ell^2(V_j(\zeta),\mu_j),
       \mathcal U_j^\zeta,A_j^\zeta,\eta_j^\zeta,
       T_{k,j}^\zeta)_{k\leq j\leq N},\\
 \mathcal F_{\rm sol}[\mathscr R]
 &:=(b,\kappa_{\rm can},C_\alpha,\kappa_\alpha)
      _{\alpha\subset\mathscr T_{k,N}(\zeta)},\\
 \mathcal F_{\rm gau}^{(m)}[\mathscr R]
 &:=(\mathscr P_j^\Sigma(\zeta),\mathscr P_j^\Pi(\zeta),
       B_j^\zeta,\kappa_j^\zeta,\delta_j^\zeta;R_m)_{k\leq j\leq N},\\
 \mathcal F_{\rm coh}^{(m)}[\mathscr R]
 &:=(\Gamma_y\otimes R_m,\partial,\Psi_\alpha,H_\bullet)
      _{y,\alpha\subset\mathscr T_{k,N}(\zeta)},\\
 \mathcal F_{\rm det}^{(m)}[\mathscr R]
 &:=(\operatorname{Det}(\operatorname{Tot}
       \mathscr D_{\bullet,\bullet;k,\ell}^{(m)}[\mathscr R]))
       _{k\leq\ell\leq N}.
\end{aligned}
\tag{C2}
\label{hmtfg:base:eq:cinco-lectores-comunes}
\]
Ici, \(\mathscr D_{\bullet,\bullet;k,\ell}^{(m)}[\mathscr R]\) est le
bicomplexe normalisé
de chemins et de mémoire des sections précédentes, restreint au sous-arbre
d’horizon \(\ell\) et réduit sur \(R_m\). Le dernier lecteur conserve la
collection ordonnée de droites déterminantales jusqu’à l’horizon \(N\) ; de cette
manière, sa troncature est une projection canonique et ne présuppose pas de
factorisation déterminantale inexistante. La loi \(\mu\) est la loi canonique
du recensement, et les ports sont fixés au moyen de la branche lexicographiquement minimale
de l’ordre APP. Par conséquent, aucun des cinq lecteurs ne reçoit de pondération
ou de port extérieurs.

Pour
\(T\in\mathfrak T:=\{{\rm sp},{\rm sol},{\rm gau},{\rm coh},{\rm det}\}\),
définissons l'espace total typé
\[
 \widetilde{\mathcal F}_{T;k,N}^{(m)}(\zeta)
 :=\coprod_{\mathscr R\in\mathfrak R_{k,N}(\zeta)}
   \{\mathcal F_T^{(m)}[\mathscr R]\}_{\mathscr R},
 \qquad
 p_T:\widetilde{\mathcal F}_{T;k,N}^{(m)}(\zeta)
      \longrightarrow\mathfrak R_{k,N}(\zeta),
\]
où \(\mathcal F_{\rm sp}^{(m)}:=\mathcal F_{\rm sp}\) et \(\mathcal F_{\rm sol}^{(m)}:=\mathcal F_{\rm sol}\) ; autrement dit, l'exposant \(m\) est inerte pour ces deux lecteurs. La projection \(p_T\) oublie uniquement la structure du lecteur et conserve son registre source. Le produit fibré correct est
\[
\begin{aligned}
 \mathcal J_{k,N}^{(m)}(\zeta):={}&
 \widetilde{\mathcal F}_{\rm sp;k,N}^{(m)}(\zeta)
 \mathop{\times}_{\mathfrak R_{k,N}(\zeta)}
 \widetilde{\mathcal F}_{\rm sol;k,N}^{(m)}(\zeta)
 \mathop{\times}_{\mathfrak R_{k,N}(\zeta)}
 \widetilde{\mathcal F}_{\rm gau;k,N}^{(m)}(\zeta)\\
 &\mathop{\times}_{\mathfrak R_{k,N}(\zeta)}
 \widetilde{\mathcal F}_{\rm coh;k,N}^{(m)}(\zeta)
 \mathop{\times}_{\mathfrak R_{k,N}(\zeta)}
 \widetilde{\mathcal F}_{\rm det;k,N}^{(m)}(\zeta).
\end{aligned}
\tag{C3}
\label{hmtfg:base:eq:constructor-correlativo}
\]
L'égalité des cinq projections oblige à conserver le même arbre, les
mêmes arêtes, la même branche produite et le même registre. Pour chaque
\(\mathscr R\), les cinq constructions déterminent l'élément
\[
 s_{k,N}^{(m)}(\mathscr R)
 :=(\mathcal F_T^{(m)}[\mathscr R])_{T\in\mathfrak T}
 \in\mathcal J_{k,N}^{(m)}(\zeta).
\]
L'existence de cet élément provient des opérations de \eqref{hmtfg:base:eq:cinco-lectores-comunes}, et non de la notation du produit fibré.

La suppression de la dernière couche induit
\[
 \operatorname{res}^{\mathfrak R}_{N+1,N}
 \bigl(\mathscr R_{k,N+1}(\zeta;x)\bigr)
 :=\mathscr R_{k,N}(\zeta;\rho_N(x)).
\]
La restriction de chaque opérateur aux générateurs retenus définit
\[
 r^m_{N+1,N}:\mathcal J_{k,N+1}^{(m)}(\zeta)
 \longrightarrow\mathcal J_{k,N}^{(m)}(\zeta).
\]
Dans le lecteur déterminantal, \(r^m_{N+1,N}\) supprime uniquement la dernière
coordonnée de la collection de droites de
\eqref{hmtfg:base:eq:cinco-lectores-comunes}. La réduction
\(R_{m+1}\twoheadrightarrow R_m\) définit, par changement de base,
\[
 q^N_{m+1,m}:\mathcal J_{k,N}^{(m+1)}(\zeta)
 \longrightarrow\mathcal J_{k,N}^{(m)}(\zeta).
\]
Sur chaque générateur, on a
\[
\begin{gathered}
 r^m_{N+1,N}q^{N+1}_{m+1,m}
 =q^N_{m+1,m}r^{m+1}_{N+1,N},\\
 \partial\Psi_\alpha=\Psi_\alpha\partial,
 \qquad L_j^\zeta\kappa_j^\zeta
       =\kappa_{j+1}^\zeta L_j^\zeta,
 \qquad L_j^\zeta\delta_j^\zeta
       =\delta_{j+1}^\zeta L_j^\zeta,
\end{gathered}
\tag{C4}
\label{hmtfg:base:eq:naturalidad-D}
\]
et les compositions des applications \(r\) et \(q\) satisfont les identités
fonctorielles. La première égalité se vérifie parce que la suppression de la dernière couche et
la réduction des coefficients agissent sur des coordonnées distinctes. Les trois autres
sont les identités de naturalité prouvées précédemment. On obtient un unique
prosystème dans \((N,m)\), et non cinq limites juxtaposées.

À la racine, on définit
\[
 \mathfrak J_{N,\bullet}^{(m)}:=\mathcal J_{0,N}^{(m)}(\ast).
\tag{C5}
\label{hmtfg:base:eq:fibras-conjuntas-horizonte}
\]
La structure discrète conjointe du continu est la limite de ces objets corrélatifs :
\[
 \boxed{
 \mathcal C_{\rm cont}^{\rm disc}
 :=\varprojlim_{\substack{N\geq1\\m\geq1}}
   \bigl(\mathfrak J_{N,\bullet}^{(m)},
         r^m_{N+1,N},q^N_{m+1,m}\bigr).}
\tag{C6}
\label{hmtfg:base:eq:continuo-discreto}
\]
C'est une définition par le prosystème engendré par l'arbre complet ; ce n'est
ni l'image ni le graphe d'une application qui retiendrait l'état d'entrée comme
première coordonnée.

Les registres possèdent une marque terminale produite. Leurs projections compatibles induisent
\[
 \Pi:\mathcal C_{\rm cont}^{\rm disc}\longrightarrow H_\infty^{\rm enr}.
\]
Réciproquement, une histoire complète détermine, par ses préfixes déjà émis, les sections compatibles \(s_{0,N}^{(m)}\). On obtient ainsi, après avoir défini le continu,
\[
 \operatorname{Gen}_{\rm cont}:H_\infty^{\rm enr}
 \longrightarrow\mathcal C_{\rm cont}^{\rm disc},
 \qquad
 \Pi\circ\operatorname{Gen}_{\rm cont}
 =\operatorname{id}_{H_\infty^{\rm enr}}.
\tag{C7}
\label{hmtfg:base:eq:seccion-historias-continuo}
\]
Cette section démontre la récupérabilité ; elle n'est pas utilisée pour définir
\(\mathcal C_{\rm cont}^{\rm disc}\).

\begin{theorem}[Émergence corrélative des cinq constructions]
\label{hmtfg:base:thm:correlacion-cinco}
Pour toute histoire survivante \(\zeta\in H_k\), tout \(N>k\), tout
\(m\geq1\) et tout \(\mathscr R\in\mathfrak R_{k,N}(\zeta)\), les cinq
constructions de \eqref{hmtfg:base:eq:cinco-lectores-comunes} sont définies sur le
même registre, et \eqref{hmtfg:base:eq:constructor-correlativo} contient
\(s_{k,N}^{(m)}(\mathscr R)\). Leurs troncatures, réductions modulaires et
transports satisfont \eqref{hmtfg:base:eq:naturalidad-D}. De plus,
\(\mathcal C_{\rm cont}^{\rm disc}\neq\varnothing\), l'application \(\Pi\)
est surjective et scindée par \(\operatorname{Gen}_{\rm cont}\). Les
cinq composantes sont des aspects inséparables du même continu, et non cinq
domaines ni cinq applications indépendantes.
\end{theorem}

\begin{proof}
La survie qui définit \(H_k\) et la loi de prolongation TPK donnent
\(\rho_k^{-1}(\zeta)\neq\varnothing\) pour chaque sommet. Comme les fibres sont
finies, le lemme de Kőnig produit \(H_\infty^{\rm enr}\neq\varnothing\).

Un registre \(\mathscr R\) étant fixé, le transport des histoires et des feuillets produit \(\mathcal F_{\rm sp}[\mathscr R]\) ; la composition affine et le défaut exact d'annulation produisent \(\mathcal F_{\rm sol}[\mathscr R]\) ; les incidences et la suite rectangulaire produisent \(\mathcal F_{\rm gau}^{(m)}[\mathscr R]\) ; le complexe de mémoire et son transport produisent \(\mathcal F_{\rm coh}^{(m)}[\mathscr R]\) ; et les droites de l'assemblage normalisé de chemins produisent \(\mathcal F_{\rm det}^{(m)}[\mathscr R]\). Chaque construction agit sur les mêmes arêtes générées, et non sur un sous-domaine supposé ni sur l'espace ambiant abstrait de \(729\) mots.

Les compatibilités affine, de transport, d’incidence et de mémoire prouvent
\eqref{hmtfg:base:eq:naturalidad-D}. Tronquer la collection déterminantale de
\eqref{hmtfg:base:eq:cinco-lectores-comunes} commute par
définition avec \(r\) ; son changement de base commute avec \(q\). Les sections
\(s_{0,N}^{(m)}\) sont compatibles, de sorte qu’elles produisent
\(\operatorname{Gen}_{\rm cont}\). L’identité
\eqref{hmtfg:base:eq:seccion-historias-continuo} prouve à la fois la non-
vacuité de la limite, la surjectivité de \(\Pi\) et la récupérabilité de chaque
histoire sans l’introduire comme coordonnée de définition.

La fibre, la relation, le nombre, l'orientation, la retenue, la mémoire, la
frontière et la route demeurent dans le registre commun ; la multisection, les
terminaux et les lecteurs linéaire et probabiliste demeurent dans les cinq
constructions corrélatives. Le passage à la limite conserve donc la clôture des
dépendances et ne disjoint pas le continu.
\end{proof}

La reconnaissance ultérieure de l'unique continu construit au moyen du langage spectral, solénoïdal, de jauge, cohomologique ou déterminantal appartient à la comparaison conventionnelle ultérieure. Elle ne sélectionne pas ses états et n'intervient pas dans sa génération APP--TRIT--TPK.
